\chapter{Introduction}
\label{ch:intro}

\section{Context}

Generative artificial intelligence reached classrooms faster than any previous
educational technology, and it arrived with an uncomfortable property: it is very good
at producing the answer a learner was supposed to produce. The immediate consequence is
that the most obvious use of a language model in education --- ``ask it and read the
reply'' --- removes exactly the cognitive work that learning consists of.

The EdNova track of the IEEE WIE ACT 5.0 competition states the problem in the opposite
direction. It asks for a companion able to act as \emph{mediator, teammate, tutor and
facilitator} so that learning becomes structured, inclusive, interactive, personalized
and effective, and it explicitly covers both individual and group settings. Read as a
specification rather than a slogan, that sentence rules out the naive design: a system
that answers for the group is a tutor at best, and a very poor mediator.

Our project starts from a narrower observation. Group study fails in ways that are
observable, repeatable and \emph{structural}:

\begin{itemize}[leftmargin=1.4em,itemsep=2pt]
  \item one participant dominates the conversation, and the group's output becomes that
        person's output;
  \item hesitant learners do not ask the ``simple'' question that would expose a shared
        misconception, because asking it in front of peers has a social cost;
  \item an unclear or wrong explanation is rarely challenged, because challenging it
        requires confidence the group does not distribute evenly;
  \item a single explanation style is imposed on everyone, so the learner who needs a
        picture, the learner who needs an analogy and the learner who needs a worked
        example all receive the same paragraph.
\end{itemize}

These are coordination and equity failures, not knowledge failures. They occur in
university tutorials, in bootcamps, in IEEE student-branch revision circles and in
informal self-learning groups alike --- which is why the same failure modes reappear
across contexts that share nothing else.

\section{Motivation and design position}

A GenAI system can address coordination failures directly, because it can observe every
message, hold state across a session, and act privately towards one learner while
speaking publicly to the group. It can also make them worse, by becoming the dominant
voice itself.

We therefore fixed a design position before writing code, and used it to arbitrate every
later decision:

\begin{keyfinding}[Design position]
The AI must not replace the circle --- it must make the circle work. Any feature that
increases the volume of AI output at the expense of peer output is, by this rule,
a regression regardless of how well it scores on answer quality.
\end{keyfinding}

Five commitments follow mechanically from that position, and they structure the rest of
this report:

\begin{enumerate}[leftmargin=1.6em,itemsep=3pt]
  \item \kw{Peers first, AI second.} Every question is offered to the circle inside a
        timed window. The AI answer is prepared during that window but withheld until it
        closes, so the group's own attempt always happens first
        (\S\ref{sec:cfas}).
  \item \kw{Protect the group from wrong peer answers.} Peer first does not mean peer
        correct. An Auditor step verifies peer claims before consensus is declared
        (\S\ref{sec:auditor}).
  \item \kw{Personalize once, then adapt.} A four-item diagnostic produces a
        learning-DNA passport that tunes level, explanation style, language register and
        hint eagerness, and a bandit updates the style weights from check-question
        outcomes (\S\ref{sec:dna}).
  \item \kw{Scaffold rather than solve.} A four-level hint ladder and Socratic
        questioning precede any direct answer (\S\ref{sec:hints}).
  \item \kw{Treat platform limits as a design budget.} Shared, GPU-less inference with a
        scarce larger-model quota and seconds-scale latency is the operating
        environment, not an incident. The routing policy, quota governor, queue and
        degradation machine exist because of it (\S\ref{ch:ai}).
\end{enumerate}

\section{Objectives}

The project pursued six objectives. Each is stated here in a form that can be checked
against the delivered system, and each is revisited with evidence in
Chapter~\ref{ch:results}.

\begin{table}[htbp]
\centering
\caption{Project objectives and the evidence that closes each one.}
\label{tab:objectives}
\footnotesize
\begin{tabularx}{\textwidth}{@{}p{0.6cm}p{5.3cm}X@{}}
\toprule
\textbf{\#} & \textbf{Objective} & \textbf{Acceptance evidence (chapter)} \\
\midrule
O1 & Structure a peer-learning session so that every learner gets a fair turn
     & Timed peer window with server-side expiry worker; peer-first rate exposed on
       \code{/metrics} (\S\ref{sec:cfas}, \S\ref{sec:obs}) \\
O2 & Detect and contain incorrect peer explanations
     & Auditor verification with poll-first fallback; seeded-misconception eval set
       (\S\ref{sec:auditor}, \S\ref{sec:eval}) \\
O3 & Personalize explanation style, level, register and language per learner
     & DNA passport, six-dimension style weights, bandit update from check outcomes
       (\S\ref{sec:dna}) \\
O4 & Keep larger-model usage a deliberate minority of calls, and prove it
     & Documented routing policy; \code{lShare} computed from the call ledger;
       measured 0.133 (\S\ref{sec:routing}, \S\ref{sec:metrics-live}) \\
O5 & Remain usable when the larger model, the quota, or the provider itself is gone
     & Four-state degradation machine with a 30-snippet offline pack
       (\S\ref{sec:degradation}) \\
O6 & Preserve learner privacy by construction, not by policy statement
     & Schema gate test that fails the build; PII screen before every model call
       (\S\ref{ch:privacy}) \\
\bottomrule
\end{tabularx}
\end{table}

\section{Scope and non-goals}

The delivered system is this final submission: a working full-stack product with a
verified backend contract, a rehearsed demo path and automated gates, not a deployed
service with real institutional users. We state the boundaries explicitly, because a
report that hides them is less useful to a jury than one that names them.

\paragraph{In scope.} Circle sessions with timed peer windows, the learning-DNA
diagnostic and passport, the hint ladder, Auditor verification, consensus polls,
revision capsules, quiz battles, the Whisper chain-repair game, S-B-U circle health,
the deterministic circle matcher and workshop board, the mentor seat, solo mode, a
facilitator dashboard, gamification with credits and XP, a four-language interface with
right-to-left support, an offline pack, and a metrics/observability surface.

\paragraph{Out of scope.} Voice interaction, real-time collaborative editing
of the board, a trained (as opposed to configured) retrieval model, federated identity
with an institutional SSO provider, and large-scale longitudinal evaluation with real
cohorts. Real-user testing with student groups in two contexts is planned for the
post-submission rehearsal and is described as planned work, not as a completed result
(\S\ref{sec:threats}).

\section{Contributions}

Beyond a working product, the report makes four engineering contributions that
generalize beyond this challenge:

\begin{enumerate}[leftmargin=1.6em,itemsep=2pt]
  \item \kw{A measured routing policy.} Rather than asserting that a two-tier split is
        efficient, \code{lShare} is derived from the persisted call ledger and exposed
        as a first-class product metric with a target, so the policy can fail review
        (\S\ref{sec:routing}).
  \item \kw{A degradation machine with semantics.} Each of the four states changes what
        the interface \emph{promises}, not just what it displays; the system never
        shows a frozen screen and never fakes streaming
        (\S\ref{sec:degradation}).
  \item \kw{Privacy as a build gate.} The absence of name, photo, gender and region
        columns is asserted by a test that fails the build if such a field appears, so
        the guarantee cannot silently erode (\S\ref{sec:privacy-gate}).
  \item \kw{A three-tier test strategy with a live smoke.} Pure-domain unit tests,
        integration tests against an ephemeral MongoDB, and a 45-check smoke against the
        real inference platform that walks the dependency order of the product
        (\S\ref{ch:testing}).
\end{enumerate}

\section{Report structure}

Chapter~\ref{ch:req} derives the requirements from the challenge statement.
Chapter~\ref{ch:arch} presents the system architecture and the two-tier model policy.
Chapter~\ref{ch:ai} details the GenAI layer: routing, retrieval, the hint ladder, the
Auditor, the Whisper game and the degradation machine.
Chapter~\ref{ch:backend} covers the backend: modules, data model, API surface and
security. Chapter~\ref{ch:frontend} covers the PWA, design system and internationalization.
Chapter~\ref{ch:privacy} covers privacy and ethics as enforced properties.
Chapter~\ref{ch:testing} describes the test strategy and its results;
Chapter~\ref{ch:results} reports measured outcomes and discusses threats to validity.
Chapter~\ref{ch:business} presents the business model, and Chapter~\ref{ch:conclusion}
concludes. The appendix lists the endpoint surface, the event taxonomy and the
reproduction procedure.
